\input{preamble.tex}

\DeclareMathOperator{\Char}{char}
\DeclareMathOperator{\im}{im}


\title{\huge{Commutative Algebra}}
\author{\LARGE{Thobias Høivik}}
\date{Fall 2026?}

\begin{document}
\maketitle

\newpage
\tableofcontents

\newpage
\section*{Intro}
These are my own personal notes in commutative algebra, 
a course I will hopefully be taking in the fall of 2026. 
Note that as these are for myself there might be gaps, inconsistencies
and whatnot which might make it hard to follow so 
keep that in mind if you are considering using my notes for 
yourself. 

\section{Rings}
Throughout the course, and these notes, a ring means 
a commutative ring with unity (meaning multiplicative 
identity).
So every ring $R$ has distinguished elements 
$1_R, 0_R \in R$ and multiplication satisfies 
$ab = ba$. 

\begin{defn}[(Commutative Unital) Ring]
    A ring is a set $A$ together with binary 
    operations $+,\cdot : A \times A \to A$ such that 
    \begin{enumerate}
        \item $(A,+)$ is an abelian group with 
            identity $0$, 
        \item multiplication is associative, 
        \item multiplication is commutative, 
        \item multiplication distributes over addition, 
        \item there exists $1 \in A$ such that 
            $1a = a$ for all $a \in A$. 
    \end{enumerate}
\end{defn}
\begin{ex}
    Examples of familiar rings include 
    \[
        \mathbb Z, \ \mathbb Q, \ \mathbb R, \ \mathbb C
    \]
    with the usual addition and multiplication. 
\end{ex}
If $A$ is a ring, then the polynomial ring $A[x]$ is 
also a ring. More generally, 
$A[x_1,\ldots,x_n]$ is the ring of polynomials in $n$ commuting 
variables with coefficients in $A$. 

\begin{defn}[Ring homomorphism]
    Let $A$ and $B$ be rings. A function 
    \[
        \phi : A \to B
    \]
    is a ring homomorphism if 
    \begin{enumerate}
        \item $\phi(1_A) = 1_B$, 
        \item $\phi(a + a') = \phi(a) + \phi(a')$, 
        \item $\phi(aa') = \phi(a) \phi(a')$. 
    \end{enumerate}
    If $\phi$ is a bijection, then it is 
    called an isomorphism. 
\end{defn}
The condition $\phi(1_A) = 1_B$ is a convention that 
one may or may not choose, but in this course it is 
not optional, and we say all homomorphisms must explicitly 
satisfy this. 

\begin{prop}
    Let $A$ be a ring. There exists a unique ring homomorphism 
    \[
        \iota : \mathbb Z \to \mathbb A.
    \]
    It is given by 
    \[
        \iota(n) = n \cdot 1_A.
    \]
\end{prop}
\begin{proof}
    Since $\iota$ is a rinh homomorphism, it must satisfy 
    \[
        \iota(1) = 1_A. 
    \]
    Then additivity forces 
    \[
        \iota(2) = \iota(1+1) = \iota(1) + \iota(1) 
        = 1_A + 1_A, 
    \]
    and similarly 
    \[
        \iota(n) = \sum_{k=1}^n 1_A = 
        \underbrace{1_A + \cdots + 1_A}_{\text{n times}} 
    \]
    for $n > 0$. Also, 
    \[
        \iota(-n) = -\iota(n).
    \]
    This there is at most one such homomorphism. 
    It is straightforward to check that the function 
    $n \mapsto n \cdot 1_A$ preserves addition, multiplication, 
    and sends $1$ to $1_A$. Hence it is a ring homomorphism. 
\end{proof}

\begin{defn}[Characteristic]
    The characteristic of a ring $A$, denoted $\Char(A)$,
    is the unique integer $n \geq 0$ such that 
    \[
        \ker(\mathbb Z \to A) = (n). 
    \]
\end{defn}
\begin{ex}
    Here are the characteristics of a few familiar 
    rings: 
    \begin{itemize}
        \item $\Char(\mathbb Z) = 0$,  
        \item $\Char(\mathbb F_p) = p$, 
        \item $\Char(\mathbb Q) = 0$. 
    \end{itemize}
\end{ex}

Intuitively, the characteristic is the smallest 
$n \in \mathbb Z^+$ such that $n \cdot 1_A = 0_A$ with 
$n := 0$ if no such $n$ exists, and in fact, the characteristic 
is usually defined as such in a first course in abstract 
algebra. 

\begin{prop}[Universal property of $\mathbb Z$-polynomials]
    Let $A$ be a ring and let $a \in A$. There exists 
    a unique ring homomorphism 
    \[
        \phi_a : \mathbb Z[x] \to A
    \]
    such that 
    \[
        \phi_a(x) = a. 
    \]
    It is given by 
    \[
        \phi_a\left(\sum_{k=0}^n c_kx^k\right) = c_0 \cdot 1_A 
        + c_1 a + \cdots + c_n a^n.
    \]
\end{prop}
\begin{proof}
    A ring homomorphism $\phi : \mathbb Z[x] \to A$ must agree 
    with the unique homomorphism $\mathbb Z \to A$ on constant 
    polynomials. 

    If $\phi(x) = a$, then multiplicativity forces 
    \[
        \phi(x^n) = a^n. 
    \]
    Additivity then forces 
    \[
        \phi(c_0 + c_1x + \cdots + c_nx^n) 
        = c_0 \cdot 1_A + c_1a + \cdots + c_na^n. 
    \]
    Thus there is at most one such homomorphism. 

    Conversely, the formula above defines a ring homomorphism, by 
        the usual rules for polynomial addition and multiplication.
\end{proof}
So homomorphism $\mathbb Z[x] \to A$ are the same thing as 
choices of an element $a \in A$. Intuitively, $\mathbb Z[x]$ is 
the "freest" ring generated by one element. 

\begin{defn}[Ideal]
    Let $A$ be a ring. An ideal $I \subseteq A$ is a subset 
    such that: 
    \begin{enumerate}
        \item $I$ is an additive subgroup of $A$, 
        \item for every $a \in A$ and every $x \in I$, we have 
            \[
                ax \in I. 
            \]
    \end{enumerate}
\end{defn}

\begin{defn}[Principal ideal]
    Let $A$ be a ring and $a \in A$. The principal ideal 
    by $a$ is 
    \[
        (a) = \{ra \bigm| r \in A\}. 
    \]
    An ideal is called principal if it is of the form 
    $(a)$ for some $a \in A$.
\end{defn}
More generally, if $a_1, \ldots, a_n \in A$, then 
\[
    (a_1,\ldots,a_n) = 
    \{r_1a_1 + \cdots + r_na_n \bigm| r_i \in A\}.
\]
This is the smallest ideal containing all of $a_1,\ldots,a_n$. 
\begin{ex}
    $(x,y) \subseteq k[x,y]$ is such an ideal. 
    \[
        (x,y) = 
        \{f(x,y)x + g(x,y)y \bigm| f,g \in k[x,y]\}. 
    \]
    It is not the whole ring. It consists precisely of 
    polynomials with zero constant term: 
    \[
        (x,y) = \{h \in k[x,y] \bigm| h(0,0) = 0\}.
    \]
\end{ex}

\begin{ex}
    $(0) = \{0\}$ is a (very uninteresting) ideal. 
    $(1) = A$ is another (very uninteresting) ideal. 
    For $n \in \mathbb Z$ we have that $(n) = n\mathbb Z$ is 
    an ideal of $\mathbb Z$. 
\end{ex}

\begin{prop}
    Let $\phi : A \to B$ be a ring homomorphism. Then 
    \[
        \ker \phi = \{a \in A \bigm| \phi(a) = 0\}
    \]
    is an ideal of $A$.
\end{prop}

\begin{proof}
    First we show that $\ker \phi$ is an additive subgroup. 
    Since $\phi(0) = 0$, we have $0 \in \ker \phi$. 
    If $x,y \in \ker \phi$, then 
    \[
        \phi(x-y) = \phi(x) - \phi(y) = 0 - 0 = 0. 
    \]
    Hence $x-y \in \ker \phi$. 
    Now let $a \in A$ and $x \in \ker \phi$. Then 
    \[
        \phi(ax) = \phi(a)\phi(x) = \phi(a) \cdot 0 = 0. 
    \]
    Thus $ax = \ker \phi$. Therefore $\ker \phi$ is an ideal. 
\end{proof}

\begin{defn}[Quotient ring]
    Let $I \subseteq A$ be an ideal. The quotient ring 
    $A/I$ is the set of additive cosets 
    \[
        a + I
    \]
    with addition and multiplication defined by 
    \[
        (a+I) + (b+I) = (a+b) + I
    \]
    and 
    \[
        (a+I)(b+I) = ab + I. 
    \]
\end{defn}
\begin{prop}
    If $I \subseteq A$ is an ideal, then $A/I$ is a ring. 
\end{prop}
\begin{proof}
    The quotient $A/I$ is already an abelian group under addition because $I$ is an additive subgroup.

    We need to check multiplication is well-defined.

    Suppose

    $$a+I=a'+I$$

    and

    $$b+I=b'+I.$$

    Then

    $$a-a'\in I,\qquad b-b'\in I.$$

    We want to prove

    $$ab+I=a'b'+I,$$

    i.e.

    $$ab-a'b'\in I.$$

    Compute:

    $$ab-a'b'=ab-a'b+a'b-a'b'=(a-a')b+a'(b-b').$$

    Since $a-a'\in I$, and $I$ absorbs multiplication by elements of $A$, we have

    $$(a-a')b\in I.$$

    Similarly,

    $$a'(b-b')\in I.$$

    Since $I$ is closed under addition,

    $$ab-a'b'\in I.$$

    Therefore multiplication is well-defined.

    The ring axioms for $A/I$ then follow from the ring axioms in $A$.
\end{proof}
The quotient map is 
\[
    \pi : A \to A/I, \ a \mapsto a + I.
\]
Its kernel is exactly $I$: 
\[
    \ker \pi = I. 
\]

\begin{thm}[Fundamental homomorphism theorem]
    Let 
    \[
        \phi : A \to B
    \]
    be a ring homomorphism. Then 
    \[
        \im \phi
    \]
    is a subring of $B$, and there is an isomorphism 
    \[
        A / \ker \phi \cong \im \phi 
    \]
    given by 
    \[
        a + \ker \phi \mapsto \phi(a). 
    \]
\end{thm}
\begin{proof}
    Define 
    \[
        \overline{\phi} : A/\ker \phi \to \im \phi
    \]
    by 
    \[
        \overline \phi(a + \ker \phi) = \phi(a). 
    \]

    We first check well-definedness.

    Suppose

    $$a+\ker\varphi=b+\ker\varphi.$$

    Then

    $$a-b\in\ker\varphi.$$

    Thus

    $$\varphi(a-b)=0.$$

    So

    $$\varphi(a)-\varphi(b)=0,$$

    and hence

    $$\varphi(a)=\varphi(b).$$

    Therefore $\overline{\varphi}$ is well-defined.

    It is a ring homomorphism because $\varphi$ is.

    It is surjective onto $\operatorname{im}\varphi$ by definition.

    Finally, its kernel is trivial. Indeed,

    $$\overline{\varphi}(a+\ker\varphi)=0$$

    if and only if

    $$\varphi(a)=0,$$

    if and only if

    $$a\in\ker\varphi,$$

    if and only if

    $$a+\ker\varphi=0+\ker\varphi.$$

    Thus $\overline{\varphi}$ is injective. 
    Hence it is an isomorphism.
\end{proof}
\begin{ex}
    \[
        \mathbb Z[x]/(x^2 + 1) \cong \mathbb Z(i) = 
        \{a + bi \bigm| a,b \in \mathbb Z\}
    \]
    by $f \mapsto f(i)$. 

    We also have 
    \[
        k[x,y]/(x) \cong k(y) 
    \]
    by $f \mapsto f(0,y)$. 

    Lastly 
    \[
        k[x,y] / (x,y) \cong k
    \]
    by $f \mapsto f(0,0)$.
\end{ex}

\begin{thm}
    Let $I \subseteq A$ be an ideal. There is a bijection 
    \[
        \{\text{ideals of} \ A/I\} 
        \leftrightarrow \{\text{ideals} \ J 
        \subseteq A \ \text{such that} \ I \subseteq J\}. 
    \]
    The maps are 
    \[
        K \subseteq A/I \mapsto \pi^{-1}(K) \subseteq A, 
    \]
    and 
    \[
        J \subseteq A, I \subseteq J \mapsto 
        J/I \subseteq A/I. 
    \]
\end{thm}

\begin{proof}
    Let 
    \[
        \pi : A \to A/I
    \]
    be the quotient map. First, if $K \subseteq A/I$ is an ideal, 
    then $\pi^{-1}(K)$ is an ideal of $A$. Since $0 \in K$, and 
    $\ker \pi = I$, we have 
    \[
        I \subseteq \pi^{-1}(K).
    \]
    Conversely, suppose $J \subseteq A$ is an ideal with 
    $I \subseteq J$. Then 
    \[
        J/I = \{j + I \bigm | j \in J\}
    \]
    is an ideal of $A/I$. Now check that these constructions are 
    inverse. If $K \subseteq A/I$, then 
    \[
        \pi(\pi^{-1}(K)) = K
    \]
    because $\pi$ is surjective. If $J \subseteq A$ contains $I$, 
    then 
    \[
        \pi^{-1}(J/I) = J. 
    \]
    Indeed, $a \in \pi^{-1}(J/I)$ if and only if 
    \[
        a + I \in J/I, 
    \]
    which happens if and only if there exists $j \in J$ such that 
    \[
        a + I = j + I.
    \]
    This means 
    \[
        a - j \in I. 
    \]
    Since $I \subseteq J$ and $j \in J$, we get 
    $a \in J$. Hence 
    \[
        \pi^{-1}(J/I) = J. 
    \]
    Therefore the two maps are inverse bijections. 
\end{proof}
\begin{ex}
    The ideals of $\mathbb Z / 12 \mathbb Z$ correspond 
    to ideals of $\mathbb Z$ containing $(12)$. 
    The ideals of $\mathbb Z$ are precisely 
    $(n)$ for $n \geq 0$. We need 
    $(12) \subseteq (n)$. This happens precisely when 
    $n \bigm| 12$. So the ideals of $\mathbb Z/12\mathbb Z$ are 
    \[
        (1), (2), (3), (4), (6), (12) = (0). 
    \]
\end{ex}

\begin{prop}
    Let $k$ be a field. The ideal 
    \[
        (x,y) \subseteq k[x,y]
    \]
    is not principal. 
\end{prop}
\begin{proof}
    Suppose for a contradiction that 
    \[
        (x,y) = (f)
    \]
    for some $f \in k[x,y]$. 
    Since 
    \[
        x \in (x,y) = (f),
    \]
    we have $f \bigm| x$. Since 
    \[
        y \in (x,y) = (f), 
    \]
    we have $f \bigm| y$. However, $x$ and $y$ have no common 
    non-unit factor in $k[x,y]$. Hence $f$ must be a unit. If 
    $f$ is a unit, then 
    \[
        (f) = (1) = k[x,y].
    \]
    But this is impossible, because 
    \[
        1 \not\in (x,y). 
    \]
    Therefore $(x,y)$ is not principal. 
\end{proof}
Thus $k[x]$ is a principal ideal domain when $k$ is a field, 
but $k[x,y]$ is not. 

\subsection*{Exercises}

\begin{prob}
    Prove that if $A$ is a ring with $1 = 0$, then 
    $A$ is the trivial ring. 
\end{prob}
\begin{proof}
    Let $a \in A$. Then 
    \begin{align*}
        a &= a \\ 
        1 \cdot a &= 0 \cdot a \\ 
        a &= 0. 
    \end{align*}
    Thus 
    \[
        A = \{0\}.
    \]
\end{proof}

\begin{prob}
    Let $A$ be a ring. Prove that there is a unique 
    homomorphism $\mathbb Z \to A$. 
\end{prob}
\begin{proof}
    Let $A$ be a ring and 
    let $\phi, \psi : \mathbb Z \to A$ be homomorphisms. 

    Then we require 
    \[
        \phi(1) = \psi(1) = 1_A. 
    \]
    Therefore, for $n > 0$, we have 
    \[
        \phi(n) = \phi\left(\sum_{k=1}^n 1\right) 
        = \sum_{k=1}^n 1_A 
        = \psi\left(\sum_{k=1}^n 1\right) 
        = \psi(n). 
    \]
    For $n < 0$ we have $n = -m$ for some 
    $m > 0$. Then 
    \[
        \phi(n) = \phi(-m) = -\phi(m) 
        = - \sum_{k=1}^n 1_A 
        = -\psi(m) = \psi(-m) = \psi(n).
    \]
    Lastly, 
    \[
        \phi(0) = 0 = \psi(0). 
    \]
    Thus $\phi$ and $\psi$ agree on all 
    of $\mathbb Z$ and hence 
    \[
        \psi = \phi. 
    \]
\end{proof}

\begin{prob}
    Identify $k[x,y] / (x,y-1)$. 
\end{prob}
\begin{proof}[Solution]
    Let $\phi : k[x,y] \to k$ by 
    \[
        \phi(f) = f(0,1). 
    \]
    Then $\phi$ is a ring homomorphism and 
    furthermore it is surjective since, for 
    $a \in k$ we have $f(x,y) := a \mapsto a$. 

    Now we show that 
    $\ker \phi = (x,y-1)$. First of all 
    $x, y-1 \in \ker \phi$ since $x(0,1) = 0$ and 
    $y-1(0,1) = 0$. Thus 
    \[
        (x,y-1) \subseteq \ker \phi.
    \]
    Now suppose $f \in \ker \phi$. Then 
    $f(0,1) = 0$ so $f(x,y) = xg(x,y) + (y-1)h(x,y) + 
    f(0,1)$. Since $f(0,1) = 0$ we have $f(x,y) = xg(x,y) + 
    (y-1)h(x,y)$ so $f(x,y) \in (x, y-1)$. Thus 
    $\ker \phi \subseteq (x,y-1)$. The two together give 
    $(x,y-1) = \ker \phi$. 
    Thus, by the fundamental theorem of ring homomorphisms 
    (or first isomorphism theorem, if you prefer), we have 
    \[
        \frac{k[x,y]}{(x,y-1)} \cong 
        k.
    \]
\end{proof}

\begin{prob}
    List all ideals of $\mathbb Z/18\mathbb Z$. 
\end{prob}
\begin{proof}[Solution]
    The ideals are 
    \begin{itemize}
        \item $(0)$, 
        \item $(1)$,
        \item $(2)$, 
        \item $(3)$, 
        \item $(6)$, 
        \item $(9)$,  
    \end{itemize}
    i.e. $(n)$ where $n \bigm| 18$. 
    Note that $(0) = (18)$ and $18$ divides $18$.
\end{proof}

\section{Special ideals and rings}

\begin{defn}
    Let $A$ be a ring. 

    An element $u \in A$ is a \textbf{unit} if there 
    exists $v \in A$ such that 
    \[
        uv = 1.
    \]
    An element $x \in A$ is a \textbf{zero-divisor} 
    if there exists a nonzero element $y \in A$ 
    such that 
    \[
        xy = 0.
    \]
    A ring $A$ is an \textbf{integral domain} if 
    \[
        A \neq 0
    \]
    and $0$ is the only zero-divisor. 

    Equivalently, $A$ is an integral domain if 
    \[
        ab = 0 \Rightarrow a = 0 \ \text{or} \ b = 0. 
    \]
    A ring $A$ is a \textbf{field} if 
    \[
        A \neq 0
    \]
    and every nonzero element of $A$ is a unit. 
\end{defn}
\begin{ex}
    $\mathbb Z$ is an integral domain, but not a field. 
    The rings 
    \[
        \mathbb Q, \ \mathbb R, \ \mathbb C
    \]
    are fields. 
    The ring $\mathbb Z/4 \mathbb Z$ is not an integral 
    domain. 
\end{ex}

\begin{prop}
    Let $A$ be a ring and $x \in A$. Then 
    \[
        x \ \text{is a unit}
    \]
    if and only if 
    \[
        (x) = (1).
    \]
\end{prop}
\begin{proof}
    Suppose first $x$ is a unit. Then $\exists y \in A$ 
    s.t. 
    \[
        xy = 1. 
    \]
    Therefore $1 \in (x)$. Since every ideal 
    containing $1$ is the whole ring, we get 
    \[
        (x) = A = (1).
    \]
    Conversely, suppose $(x) = (1)$. 
    Then 
    \[
        1 \in (x). 
    \]
    By definition of the principal ideal $(x)$, there exists 
    $y \in A$ such that 
    \[
        1 = yx. 
    \]
    Since $A$ is commutative, 
    \[
        xy = 1.
    \]
    Therefore $x$ is a unit.
\end{proof}

\begin{thm}
    Let $A$ be a ring. Then $A$ is a field if and only if 
    $A$ has exactly two ideals: 
    \[
        (0) \neq (1). 
    \]
\end{thm}
\begin{proof}
    Suppose $A$ is a field. 

    Let $I \subseteq A$ be a nonzero ideal. Choose 
    \[
        0 \neq x \in I. 
    \]
    Since $A$ is a field, $x$ is a unit. Then 
    \[
        (x) = (1). 
    \]
    But 
    \[
        (x) \subseteq I. 
    \]
    Therefore 
    \[
        (1) = A \subseteq I
    \]
    so 
    \[
        I = (1). 
    \]
    So the only ideals are 
    \[
        (0), \ (1). 
    \]

    Conversely, suppose $A$ has two ideals 
    \[
        (0) \neq (1). 
    \]
    Let $0 \neq x \in A$. Then 
    \[
        (x) \neq (0). 
    \]
    Since the only nonzero ideal is $(1)$, we have 
    \[
        (x) = (1). 
    \]
    Then $x$ is a unit and thus since every nonzero element of 
    $A$ is a unit, $A$ is a field. 
\end{proof}

\begin{defn}[Prime ideal]
    Let $A$ be a ring and let $I\subseteq A$ be an ideal. 
    We say $I$ is a \textbf{prime ideal} if

    $$I\neq(1)$$

    and whenever

    $$ab\in I,$$

    we have

    $$a\in I\quad\text{or}\quad b\in I.$$

    Equivalently,

    $$a,b\notin I\implies ab\notin I.$$

    The set of all prime ideals of $A$ is denoted

    $$\operatorname{Spec}(A).$$
\end{defn}

\begin{ex}
    The prime ideals of $\mathbb Z$ are 
    $(0)$ and $(p)$ where $p$ is a prime number. 
\end{ex}

\begin{defn}[Maximal ideal]
    Let $A$ be a ring and let $I\subseteq A$ be an ideal. 
    We say $I$ is a \textbf{maximal ideal} if

    $$I\neq(1)$$

    and the only ideals $J$ satisfying

    $$I\subseteq J\subseteq A$$

    are

    $$J=I$$

    and

    $$J=A.$$

    Equivalently, there is no ideal strictly between $I$ and $A$.

    The set of maximal ideals of $A$ is denoted

    $$\operatorname{Specm}(A).$$
\end{defn}
\begin{ex}
    The maximal ideals of $\mathbb Z$ are $(p)$. 
\end{ex} 

\begin{thm}
    Let $A$ be a ring and let $I \subseteq A$ be an ideal. 
    Then 
    \[
        I \ \text{is prime}
    \]
    if and only if 
    \[
        A/I \ \text{is an integral domain}.
    \]
    Also, 
    \[
        I \ \text{is maximal}
    \]
    if and only if 
    \[
        A/I \ \text{is a field}. 
    \]
\end{thm}

\subsection*{Exercises}

\begin{prob}
    Decide whether $(2) \subseteq \mathbb Z/6\mathbb Z$ 
    is prime, maximal, both, or neither. 
\end{prob}
\begin{proof}[Solution]
    It is maximal. 
    Let $\phi : \mathbb Z/6\mathbb Z \to \mathbb Z/2\mathbb Z$ 
    by 
    \[
        \phi(a) = a \mod 2. 
    \]
    Then this map is surjective and 
    $\ker \phi = (2)$. By the first isomorphism theorem we 
    have 
    \[
        \mathbb Z_6/\ker \phi \cong \mathbb Z_2
    \]
    which is a field and thus $\ker \phi$ is maximal.  
\end{proof}

\begin{prob}
    Decide whether $(x-1) \subseteq k[x]$ is prime or maximal. 
\end{prob}
\begin{proof}[Solution]
    Let $\phi : k[x] \to k$ by 
    \[
        \phi(f) = f(1). 
    \]
    Then $\ker \phi = (x-1)$, the map is surjective so 
    it follows as from before that $(x-1)$ is maximal. 
\end{proof}

\section{Krull's theorem, PIDs/UFDs, local rings}

\begin{thm}[Krull's theorem]
    Let $A$ be a nonzero ring. Then $A$ has a maximal 
    ideal. 

    Equivalently: 
    \[
        A \neq 0 
        \Rightarrow \operatorname{Specm}(A) \neq \emptyset. 
    \]
\end{thm}
This is not at all obvious for arbitrary rings, because 
rings may have infinitely many ideals and no obvious "largest" 
proper ideal. To prove Krull's theorem we will need 
the following set-theoretic principle. 

\begin{lem}[Zorn's lemma]
    Let $S$ be a partially ordered set. Suppose 
    every chain in $S$ has an upper bound in $S$. 
    Then $S$ has a maximal element. 
\end{lem}

Here: 
\begin{itemize}
    \item a \textbf{chain} is a subset where any 
        two elements are comparable, 
    \item an \textbf{upper bound} for a chian 
        $C \subseteq S$ is an element $u \in S$ such 
        that $c \leq u$ for all $u \in C$, 
    \item a \textbf{maximal element} is an element 
        $m \in S$ such that there is no $s \in S$ 
        with $m < s$. 
\end{itemize}

\begin{proof}[Proof of Krull's theorem]
    Let $A \neq 0$. Consider the set 
    \[
        S = \{I \subseteq A \bigm| I \ \text{is an ideal and} \ 
        I \neq A\}.
    \]
    So $S$ is the set of all proper ideals of $A$. 
    We partially order $S$ by inclusion: 
    \[
        I \leq J \Leftrightarrow I \subseteq J. 
    \]
    We want to apply Zorn's lemma. 

    Let $\{I_\lambda\}_{\lambda \in \Lambda}$ be a chain in 
    $S$. Define 
    \[
        I = \bigcup_{\lambda \in \Lambda} I_\lambda.
    \]
    We claim $I \in S$, meaning $I$ is a proper ideal. 

    First, $I$ is an ideal. 

    Let $a,b \in I$. Then there exists $\lambda, \mu \in \Lambda$ 
    such that 
    \[
        a \in I_\lambda, \ b \in I_\mu. 
    \]
    Since $I_\mu$ is an ideal, 
    \[
        a - b \in I_\mu \subseteq I. 
    \]
    Thus $I$ is an additive subgroup. 
    Now let $r \in A$ and $a \in I$. Then $a \in I_\lambda$ for 
    some $\lambda$, so 
    \[
        ra \in I_\lambda \subseteq I. 
    \]
    Thus $I$ absorbs multiplication from $A$ and is therefore 
    an ideal. Second $I \neq A$. 
    If $I = A$, then in particular 
    \[
        1 \in I.
    \]
    So $1 \in I_\lambda$ for some $\lambda$, but then 
    \[
        I_\lambda = A,
    \]
    contradicting $I_\lambda \in S$. Hence $I \neq A$. 

    So $I \in S$, and $I$ is an upper bound for the chain. 
    Therefore every chain in $S$ has an upper bound. 
    By Zorn's lemma, $S$ has a maximal element $m$. 
    This $m$ is a proper ideal of $A$, and no proper ideal 
    strictly contains it. Hence $m$ is a maximal ideal. 

    Thus every nonzero ring has a maximal ideal. 
\end{proof}

\begin{cor}
    Let $I \subseteq A$ be a proper ideal. 
    Then there exists a maximal ideal $m \subseteq A$ such 
    that 
    \[
        I \subseteq m. 
    \]
\end{cor}
\begin{proof}
    Since $I \neq A$, the quotient 
    \[
        A/I
    \]
    is nonzero. By Krull's theorem, $A/I$ has a maximal ideal 
    $\mathfrak m$. 
    Let 
    \[
        \pi : A \to A/I   
    \]
    be the quotient map. Then 
    \[
        m = \pi^{-1}(\mathfrak m)
    \]
    is a maximal ideal of $A$ containing $I$, by 
    the correspondence theorem. 
\end{proof}
\begin{cor}
    Let $A$ be a ring and $x \in A$. Then 
    \[
        x \ \text{is a unit}
    \]
    if and only if 
    \[
        x \not\in m
    \]
    for every maximal ideal $m \subseteq A$. 

    Equivalently: 
    \[
        x \ \text{is not a unit} 
        \Leftrightarrow x \in m \ \text{for some maximal ideal} \ 
        m. 
    \]
\end{cor}
\begin{proof}
    We have: 
    \[
        x \ \text{is not a unit} \Leftrightarrow 
        (x) \neq A. 
    \]
    By the previous Corollary, 
    \[
        (x) \neq A \Leftrightarrow (x) \subseteq m
    \]
    for some maximal ideal $m$. 
    But 
    \[
        (x) \subseteq m \Leftrightarrow x \in m. 
    \]
    Therefore 
    \[
        x \ \text{is not a unit} 
        \Leftrightarrow x \in m
    \]
    for some maximal ideal $m$. 
\end{proof}

\begin{defn}[PID]
    An integral domain $A$ is a 
    \textbf{principal ideal domain}, or PID, if every ideal 
    of $A$ is principal. 
\end{defn}
\begin{ex}
    Some examples of PIDs include 
    \begin{itemize}
        \item $k[x]$ for any field $k$, 
        \item any field $k$, 
        \item $\mathbb Z$. 
    \end{itemize}
\end{ex}

\begin{defn}
    Let $A$ be an integral domain. Let $f \in A$ be neither 
    $0$ nor a unit. We say $f$ is \textbf{irreducible} 
    if whenever 
    \[
        f = gh, 
    \]
    then either $g$ is a unit or $h$ is a unit, 
    meaning $f$ cannot be factored nontrivially. 
\end{defn}
\begin{defn}[UFD]
    An integral domain $A$ is a \textbf{unique factorization domain},
    or UFD, if every nonzero nonunit $f \in A$ can be written as 
    a finite product of irreducibles, 
    \[
        f = p_1 \cdots p_n,
    \]
    and this factorization is unique up to reordering and  
    multiplication by units. 
\end{defn}
\begin{ex}
    Some examples include: 
    \begin{itemize}
        \item $\mathbb Z$, 
        \item $k[x]$ for any field $k$, 
        \item $k[x_1,\ldots,x_n]$ for any field $k$. 
    \end{itemize}
\end{ex}

We remark that PID $\Rightarrow$ UFD and that if 
$A$ is a UFD then $A[x]$ is a UFD. 

\begin{defn}[Local ring]
    A ring $A$ is \textbf{local} if it has exactly 
    one maximal ideal. 
    That maximal ideal is usually denoted 
    \[
        \mathfrak m. 
    \]
    So a local ring is a pair $(A, \mathfrak m)$, 
    where $\mathfrak m$ is the unique maximal ideal. 
\end{defn}

\begin{ex}
    Every field $k$ is local since the only maximal ideal 
    is $(0)$. 

    $\mathbb Z$ is a nonexample as it has maximal ideals 
    $(p)$ for every prime $p$. 
\end{ex}

\begin{prop}
    Let $A$ be a local ring with unique maximal ideal 
    $\mathfrak m$. Then the set of units in $A$ is 
    \[
        A \setminus \mathfrak m. 
    \]
    Equivalently: 
    \[
        x \in \mathfrak m 
        \Leftrightarrow x \ \text{is not a unit}. 
    \]
\end{prop}
\begin{proof}
    By the corollary to Krull's theorem, an element $x \in A$ 
    is a unit if and only if it is contained in no maximal ideal. 
    But $A$ has exactly one maximal ideal, namely 
    $\mathfrak m$. 
    Therefore 
    \[
        x \ \text{is a unit} \Leftrightarrow 
        x \not\in \mathfrak m. 
    \]
    So the units are exactly 
    \[
        A \setminus \mathfrak m. 
    \]
\end{proof}

\begin{prop}
    Let $A$ be a ring. Suppose the set of nonunits 
    \[
        \mathfrak m = \{x \in A \bigm| x \ \text{is not 
        unit}\}
    \]
    is an ideal. 
    Then $A$ is local, and $\mathfrak m$ is its unique 
    maximal ideal. 
\end{prop}
\begin{proof}
    Let $I \subseteq A$ be a proper ideal. 
    We claim 
    \[
        I \subseteq \mathfrak m. 
    \]
    Indeed, if $x \in I$ were a unit, then 
    \[
        1 \in I,
    \]
    because 
    \[
        x^{-1}x = 1. 
    \]
    Then $I = A$, contradicting that $I$ is proper. 
    So every element of $I$ is a nonunit, hence 
    \[
        I \subseteq \mathfrak m. 
    \]
    Thus every proper ideal of $A$ is contained in 
    $\mathfrak m$. 
    Since $\mathfrak m$ is itself a proper ideal, 
    it is the unique maximal ideal. Therefore $A$ is local. 
\end{proof}

\subsection*{Exercises}

\begin{prob}
    \ 
    \begin{enumerate}[label=(\alph*)]
        \item Let $A$ be a ring, and let 
            $f = a_nx^n + \cdots + a_0 \in A[x]$. Prove 
            that if $f$ is a unit of $A[x]$, then 
            $a_0$ is a unit of $A$.
        \item Show that if $a$ is a unit of $A$, then the 
            constant polynomial $a \in A[x]$ is a unit of 
            $A[x]$. 
        \item Show that if $A$ is an integral domain, 
            then the units of $A[x]$ are exactly the constant 
            polynomials $a$ with $a$ a unit of $A$. 
    \end{enumerate}
\end{prob}
\begin{proof}[Proof of (a)]
    Suppose 
    \[
        f = \sum_{k=0}^n a_kx^k
    \]
    is a unit of $A[x]$. Then there exists 
    \[
        g = \sum_{k=0}^m b_kx^k
    \]
    such that 
    \[
        fg = 1. 
    \]
    Let $\phi : A[x] \to A$ by 
    \[
        \phi(f) = f(0).
    \]
    Then $\phi$ is a ring homomorphism, and thus 
    \[
        1_A = \phi(1_{A[x]}) 
        = \phi(fg) = a_0 b_0 
    \]
    so $a_0$ is a unit in $A$. 
\end{proof}
\begin{proof}[Proof of (b)]
    Suppose $a \in A$ is a unit. Then there 
    exists $b \in A$ so that 
    \[
        ab = 1_A. 
    \]
    Then $a,b \in A[x]$ as constant polynomials. 
    Once again 
    \[
        1_{A[x]} = 1_A = 
        ab = \phi(a)\phi(b) 
        = \phi(ab). 
    \]
\end{proof}
\begin{proof}
    From the above we have that 
    units in $A$ as constant polynomials 
    are units in $A[x]$.  

    Now let $f \in A[x]$ be a unit. Then 
    there is $g \in A[x]$ so 
    \[
        fg = 1. 
    \]
    Then since $\deg(g\cdot h) = \deg(g) + \deg(h)$ we have 
    \[
        \deg(f) + \deg(g) = \deg(f\cdot g) = \deg(1) = 0. 
    \]
    Then $\deg(f) = \deg(g) = 0$ so they are constant polynomials. 

    We require $A$ to be an integral domain because otherwise 
    we could take $A = \mathbb Z_4$ and $f = 2x + 1$. 
    Then $f^2 = 1$. 
\end{proof}

\begin{prob}
    Prove that a ring $A$ is a field if and only if 
    the following holds: $A \neq 0$ and every homomorphism 
    from $A$ to a non-zero ring is injective. 
\end{prob}
\begin{proof}
    Let $A$ and $B \neq 0$ be rings. 

    First assume $A$ is a field. 
    Then $\ker \phi$ is $A$ or $(0)$ for every homomorphism 
    $\phi : A \to B$. If $\ker \phi = A$ then 
    $\psi(x) = 0_B$ for all $x \in A$ and in particular 
    $1_B = \phi(1_A) = 0_B$ so $B$ is the trivial ring, but 
    this contradicts our assumption. Thus 
    $\ker \phi = (0)$, meaning $\phi$ is injective. Furtheremore, 
    by definition, $A \neq 0$.
    
    Now assume $A \neq 0$ and every homomorphism 
    $\phi : A \to B$ has $\ker \phi = \{0_A\}$ when $B$ is nonzero. 

    Let $I \subseteq A$ is some ideal. Consider the canonical 
    quotient homomorphism 
    \[
        \pi : A \to A/I \ \text{defined by} \ 
        \pi(a) = a+I. 
    \]
    Now either $A/I$ is nonzero or not. If it is the zero ring 
    then $\ker \pi = A$. 
    Otherwise if $A/I$ is nonzero, then by assumption 
    $\ker \pi = (0)$. By construction, 
    \[
        \ker \pi = I
    \]
    so every ideal $I$ of $A$ is either $(0)$ or $A$. Hence 
    $A$ is a field. 
\end{proof}

\begin{prob}
    Let $A$ be a ring, and let $x \in A$. Consider 
    the map $m_x : A \to A$ given by $m_x(y) = xy$ for 
    some $x \in A$. 
    \begin{enumerate}[label=(\alph*)]
        \item Show that $m_x$ is a ring homomorphism if 
            and only if $x = 1$. 
        \item Show that $x$ is a $0$-divisor if and only if 
            $m_x$ is non-injective. 
        \item Show that $x$ is a unit if and only if $m_x$ 
            is surjective. 
    \end{enumerate}
\end{prob}
\begin{proof}[Proof of (a)]
    Suppose $x = 1$. Then $m_x : A \to A$ is the 
    identity map $\operatorname{id} : A \to A$, 
    \[
        \operatorname{id}(y) = y,  
    \]
    which is a homomorphism. 

    Now suppose $m_x$ is a ring homomorphism and suppose 
    for a contradiction that $x \neq 1$. 
    Then 
    \[
        m_x(1) = x \neq 1 
    \]
    so $m_x$ does not preserve the multiplicative 
    identity and thus fails to be a homomorphism. 
\end{proof}

\begin{proof}[Proof of (b)]
    Suppose $x$ is a $0$-divisor, meaning $x \neq 0$ and there  
    exists $\xi \in A\setminus \{0\}$ so that 
    \[
        x \cdot \xi = 0.
    \]
    Then $\xi \in \ker m_x$ and $\xi \neq 0$ so 
    $m_x$ cannot be injective. 

    Suppose $m_x$ is not injective. 
    So $|\ker m_x| > 1$. Then there exists 
    $\xi \in \ker m_x \setminus \{0\}$. 
    Then 
    \begin{align*}
        m_x(\xi) &= m_x(0) \\ 
        x \cdot \xi = x \cdot 0 = 0
    \end{align*}
    so $x$ is a $0$-divisor.
\end{proof}
\begin{proof}[Proof of (c)]
    Suppose $x$ is a unit. Then there exists 
    $x^{-1} \in A$ so that 
    $xx^{-1} = 1$. 
    Then, for any $a \in A$ we see that 
    \[
        xx^{-1}a \overset{m_x}{\mapsto} a
    \]
    so $m_x$ is surjective. 

    Now suppose $m_x$ is surjective. Then, for 
    any $a \in A$ there exists $b \in A$ so that 
    \[
        m_x(b) = xb = a. 
    \]
    In particular this holds 
    for $1 \in A$ so there exists $x^{-1} \in A$ so that 
    \[
        m_x(x^{-1}) = xx^{-1} = 1. 
    \]
\end{proof}

\begin{prob}
    Let $A$ be a ring, and let $\{P_i\}_{i\in I}$ be a 
    non-empty chain of prime ideals (where "chain" ,eans 
    that for every two $i,j \in I$ either $P_i \subseteq P_j$ 
    or $P_j \subseteq P_i$). Prove that 
    $\bigcup_{i\in I}P_i$ and $\bigcap_{i\in I}P_i$
    are prime ideals. 
\end{prob}
\begin{proof}
    Let 
    \[
        J = \bigcup_{i \in I} P_i
    \]
    and 
    \[
        K = \bigcap_{i \in I} P_i.
    \]

    We begin with $J$. 
    Since the chain is nonempty, pick any $P_i$. Since 
    $P_i$ is an ideal, $0 \in P_i$ so $0 \in J \neq \emptyset$. 
    Let $x \in J$ and $r \in A$. Since $x \in J$, 
    $x \in P_i$ for some specific index $i$. Since $P_i$ is 
    an ideal we have $rx \in P_i$ so $rx \in J$. 
    Let $x,y \in J$. Then $x \in P_i$ and $y \in P_j$ for 
    some specific indicies $i,j\in I$. The chain-condition 
    gives us that $P_i \subseteq P_j$ or vice-versa. Assume 
    wlog the former holds. Then 
    $x,y \in P_j$ so $x+y \in P_j \subseteq J$. Thus 
    $J$ is closed under addition. 

    Now since every $P_i$ is prime we have $1 \not\in P_i$ for 
    every $i \in I$ so $1 \not \in J$. Thus $J \neq A$ and is 
    therefore a proper ideal. Let $a,b \in A$ such 
    that $ab \in J$. Then there is some $P_i$ so that 
    $ab \in P_i$. Since $P_i$ is prime we have 
    $a \in P_i$ or $b \in P_i$. 

    Now we do $K$. 
    Similar as before $0 \in P_i$ for all $i \in I$ so 
    $0 \in K \neq \emptyset$. 
    Let $x \in J$ and $r \in A$. Since $x \in J$ we have that 
    $x \in P_i$ for every $i$ so $ax \in P_i$ for every 
    $i$ hence $ax \in K$. Let $x,y \in K$. Then 
    $x,y \in P_i$ for every $i$ hence 
    $x +y \in P_i$ for every $i$ so $x+ y \in K$. 
    Thus $K$ is closed under addition. 

    Now supose $ab \in K$. Then $ab \in P_i$ for every 
    $i$ so $a \in P_i$ or $b \in P_i$ for every $i$. 
    Thus $a \in K$ or $b \in K$. 
\end{proof}

\begin{prob}
    Let $X$ be a topological space and $x \in X$. 
    We consider pairs $(U,f)$, where $U \subseteq X$ is an open 
    set with $x \in U$ and $f : U \to \mathbb R$ is a 
    continuous function. 

    Two pairs $(U,f)$ and $(V,g)$ are called equivalent if 
    there exists an open neighborhood $W$ of $x$ with 
    $W \subseteq U \cap V$ such that 
    $f\bigm|_W = g\bigm|_W$. An equivalence class is called 
    a germ of a function at $x$. We denote by $\mathcal O_{X,x}$ 
    the set of all such germs. 

    \begin{enumerate}[label=(\alph*)]
        \item Show that addition and multiplication of germs 
            are well defined by 
            \[
                (U,f) + (V,g) 
                := (U \cap V, f\bigm|_{U\cap V} 
                + g\bigm|_{U\cap V}), 
            \]
            \[
                (U,f) + (V,g) 
                := (U \cap V, f\bigm|_{U\cap V} 
                g\bigm|_{U\cap V}). 
            \]
            Verify that $\mathcal O_{X,x}$ is a ring 
            with these operations, and describe the 
            elements $0,1 \in \mathcal O_{X,x}$ as germs. 
        \item Show that a germ $(U,f)$ is a unit in 
            $\mathcal O(X,x)$ if and only if $f(x) \neq 0$. 
        \item Show that $\mathcal O_{X,x}$ is a local ring, 
            with maximal ideal 
            \[
                \mathfrak m = \{(U,f)\in \mathcal O_{X,x} 
                \bigm| f(x) = 0\}, 
            \]
            and residue field 
            \[
                \mathcal O_{X,x}/\mathfrak m \cong \mathbb R.
            \]
    \end{enumerate}
\end{prob}
\begin{proof}[Proof of (a)]
    Suppose $(U_1,f_1) = (U_2, f_2)$ and 
    $(V_1, g_1) = (V_2, g_2)$. By definition of the equivalence 
    relation 
    \begin{itemize}
        \item there exists an open neighborhood 
            $W_1 \subseteq U_1\cap U_2$ of $x$ such that 
            $f_1\bigm|_{W_1} = f_2\bigm|_{W_1}$, and 
        \item there exists an open neighborhood 
            $W_2 \subseteq V_1\cap V_2$ of $x$ such that 
            $g_1\bigm|_{W_2} = g_2\bigm|_{W_2}$.
    \end{itemize}
    Define $W = W_1 \cap W_2$. Since open neighborhoods 
    are closed under finite intersections we have that 
    $W$ is an open neighborhood of $x$. Furthermore, 
    $W \subseteq (U_1 \cap v_1) \cap (U_2 \cap V_2)$. 
    Restricting the sum to $W$ we get 
    \[
        (f_1 + g_2) \bigm|_W 
        = f_1\bigm|_W + g_1\bigm|_W 
        + f_2\bigm|_W + g_2\bigm|_W 
        = 
        (f_2 + g_2) \bigm|_W. 
    \]
    Thus addition is well-defined and this approach works  
    for multiplication as well. 

    The ring properties of associativity, commutativity and 
    distributivity are inherited from the pointwise operations 
    of $\mathbb R$. 

    The additive itentity is $(X,\textbf{0})$ where 
    $\textbf{0} : X \to \mathbb R$ is the constant 
    function $0$. 

    The multiplicative identity is 
    $(X, \textbf{1})$ where $\textbf{1} : X \to \mathbb R$ is 
    the constant $1$. 
\end{proof}
\begin{proof}[Proof of (b)]
    Suppose $(U,f)$ is a unit. Then there exists 
    $(V,g)$ such that $(U,f)\cdot (V,g) = (X,\textbf{1})$. 
    Then there is an open neighborhood $W \subseteq U \cap V$ 
    of $x$ where $f(k)g(k) = 1$ for all $k \in W$. 
    Evaluated at $x$ we get 
    \[
        f(x)g(x) = 1. 
    \]
    Since $f(x),g(x) \in \mathcal R$ we cannot have 
    $f(x) = 0$ as $0^{-1}$ does not exist.  

    Suppose $f(x) \neq 0$. Since $f: U \to \mathbb R$ is 
    continuous and $\{0\}^c = \mathbb R \setminus \{0\}$ is 
    an open subset of $\mathbb R$, the preimage 
    $V = f^{-1}(\{0\}^c)$ is open in $U$ and thus open 
    in $X$. 

    Since $f(x) \neq 0$ we have $x \in V$ so we can define 
    a continuous $g : V \to \mathbb R$ by 
    \[
        g(k) = \frac{1}{f(k)} \ \text{for all} \ k \in V.
    \]
    Then, on the open neighborhood $V$, we have 
    $f\bigm|_V \cdot g = 1\bigm|_V$. Thus 
    $[U,f]$ is a unit. 
\end{proof}


\end{document}
